%----------------------------------------------------------------------
%%% LEARNING
%----------------------------------------------------------------------
% !TEX root = ./Main.tex


We now proceed to describe our general model for episodic learning in stochastic games.
%In tune with the episodic framework described in the previous section, we will likewise
To that end, we will consider a framework where agents follow a specific policy $\curr$ within each episode, and update it from one episode to the next with the objective of increasing their individual rewards.
Formally, our approach will adhere to the following inter-episode sequence of events:
\begin{enumerate}
\item
At the beginning of each episode $\run=\running$, every agent $\play\in\players$ chooses a policy $\policy_{\play,\run} \in \policies_{\play}$.
\item
Within the $\run$-th episode, each player executes their chosen policy $\policy_{\play,\run}$, inducing in this way an intra-episode trajectory of play $\traj_{\run} = (\rstate^{(\run)}_{\time}, \pure^{(\run)}_{\time}, \reward^{(\run)}_{\time})_{\time \leq \stopTime(\traj_{\run})}$.
\item
Once the episode terminates, agents update their policies and the process repeats.
\end{enumerate}

In terms of feedback, we will treat several models, depending on what type of information is available to the agents during play.
More precisely, we will focus on the generic \ac{PG} template
\begin{equation}
\label{eq:PG}
\tag{PG}
\next
	= \proj_{\policies}(\curr + \curr[\step] \curr[\signal])
\end{equation}
where:
\begin{enumerate}
%[(\itshape i\hspace*{1pt}\upshape)]
\item
$\curr = (\state_{\play,\run})_{\play\in\players} \in \policies$ denotes the player's policy profile at each episode $\run=\running$
\item
$\curr[\signal] = (\signal_{\play,\run})_{\play\in\players} \in \dspace$ is an estimate for the agents' inidividual policy gradients.
\item
$\proj_{\policies}\from \dspace \to \policies$ denotes the Euclidean projection to the agents' policy space $\policies$.
\item
$\curr[\step] > 0$ is the method's step-size,
for which we will assume throughout that $\sum_{\run} \curr[\step] = \infty$;
typically, \eqref{eq:PG} is run with a step-size of the form $\curr[\step] = \step/(\run + \runoff)^{\pexp}$ for some $\step>0$, $\runoff\geq0$ and $\pexp\geq0$.
\end{enumerate}

Regarding the gradient signal $\curr[\signal]$, we will decompose it as
\begin{equation}
\label{eq:signal}
\curr[\signal]
	= \vecfield(\curr) + \curr[\noise] + \curr[\bias]
\end{equation}
where
\begin{equation}
\label{eq:bias-noise}
\curr[\noise]
	= \curr[\signal] - \exof{\curr[\signal] \given \curr[\filter]}
	\quad
	\text{and}
	\quad
\curr[\bias]
	= \exof{\curr[\signal] \given \curr[\filter]} - \vecfield(\curr).
\end{equation}
In the above, we treat $\curr$, $\run=\running$, as a stochastic process on some complete probability space $\probspace$, and we write $\curr[\filter] \defeq \filter(\init,\dotsc,\curr) \subseteq \filter$ for the history (adapted filtration) of $\curr$ up to \textendash\ and including \textendash\ stage $\run$.
By definition, $\exof{\curr[\noise] \given \curr[\filter]} = 0$ and $\curr[\bias]$ is $\curr[\filter]$-measurable, so $\curr[\noise]$ can be intepreted as a random, zero-mean error relative to $\vecfield(\curr)$, whereas $\curr[\bias]$ captures all systematic (non-zero-mean) errors.
To make this precise, we will further assume that $\curr[\bias]$ and  $\curr[\noise]$  are bounded as
\begin{equation}
\label{eq:errorbounds}
\exof{\dnorm{\curr[\bias]} \given \curr[\filter]}
	\leq \curr[\bbound]
	\qquad
	\text{ and }
	\qquad
\exof{\dnorm{\curr[\noise]}^{2} \given \curr[\filter]}
	\leq \curr[\sdev]^{2}
%	\qquad
%	\text{and}
%	\qquad
%\dnorm{\curr[\signal]}
%	\leq \curr[\sbound]
\end{equation}
where the sequences $\curr[\bbound]$ and $\curr[\sdev]$ %and $\curr[\sbound]$ 
, $\run=\running$, are to be construed as deterministic upper bounds on the bias, fluctuations, and magnitude of the gradient signal $\curr[\signal]$.

Depending on these bounds, a gradient signal with $\curr[\bbound]=0$ will be called \emph{unbiased},
%and a signal with $\lim_{\run\to\infty}\curr[\bbound] = 0$ will be called \emph{asymptotically unbiased};
%finally,
and an unbiased signal with $\curr[\sdev] = 0$  will be called \emph{perfect}.
More generally, we will assume that the above statistics are bounded as
\begin{equation}
\label{eq:schedule}
\curr[\bbound]
	= \bigoh(1/\run^{\bexp})
	\qquad
	\text{and}
	\qquad
\curr[\sdev]
	= \bigoh(\run^{\sexp})
\end{equation}
for some $\bexp,\sexp>0$ which depend on the specific model under consideration.
For concreteness, we describe below three basic models that adhere to the above template for $\curr[\signal]$ in order of decreasing information requirements:


\begin{model}
[Full gradient information]
\label{mod:full}
The first model we will consider assumes that agents observe their \emph{full policy gradients}, \ie
\begin{equation}
\label{eq:signal-full}
\curr[\signal]
	= \gvalue(\curr)
\end{equation}
implying in particular that $\curr[\noise] = \curr[\bias] = 0$.
%(so we can formally take $\bexp = \infty$ and $\sexp=-\infty$).
This model is fully deterministic across episodes (though intra-episode play remains stochastic).
In particular, it tacitly assumes that agents know the game (and can observe their opponents' policies) so as to calculate the full gradients of their individual value functions $\rvalue_{\play,\stateDist}$, \cf \cite{zhang2021gradient,agarwal2021theory,LOPP22} and references therein.
\endenv
\end{model}


\begin{model}
[Learning with stochastic gradients]
\label{mod:stoch}
A relaxation of the above model which is particularly relevant for applications to deep reinforcement learning concerns the case where the player have access to stochastic policy gradients \cite{zhang2020global}, \ie unbiased gradient estimates of the form
\begin{equation}
\label{eq:signal-stoch}
\curr[\signal]
	= \gvalue(\curr)
	+ \curr[\noise]
\end{equation}
with $\exof{\curr[\noise] \given \curr[\filter]} = 0$ (so we can formally take $\bexp = \infty$ and $\sexp = 0$ in \cref{eq:schedule}
 above).
%This case is considered in \cite{zhang2020global}.
\endenv
\end{model}


%----------------------------------------------------------------------
%% Reinforce/Greedy algo begins here

\begin{figure*}[t]
\begin{minipage}{0.49\textwidth}
\begin{algorithm}[H]
\caption{\reinforce}
\label{alg:reinforce}
\begin{algorithmic}[1]
	\State 
	\textbf{Input:} $\estpolicy \in \policies, \traj = (\rstate_{\time}, \pure_{\time}, \reward_{\time})_{\time \le \stopTime(\traj)} \in\mathcal{T}$\vskip 3pt
	\For{$\play = 1,\hdots, \nPlayers$}\vskip 3pt
	\State 
		$\rewards_\play\parens{\traj} \leftarrow \sum_{\time=\tstart}^{\stopTime\parens{\traj}}\reward_{\play,\time}$\vskip 3pt
	\State 
		$\logtrick_\play\parens{\traj} \leftarrow \sum_{\time=\tstart}^{\stopTime\parens{\traj}} \nabla_{\play} \parens*{\log\estpolicy_\play\parens{\action_{\play,\time}|\rstate_\time}}$\vskip 3pt
	\State $\logvalue_\play \leftarrow \rewards_\play\parens{\traj}\cdot\logtrick_\play\parens{\traj}$\vskip 3pt
	\EndFor\vskip 3pt
	\State
	\Return $\braces{\logvalue_\play}_{\play\in\players}$
	\end{algorithmic}
\end{algorithm}
\end{minipage}
\hfill
\begin{minipage}{0.49\textwidth}
\begin{algorithm}[H]
\caption{\greedy}
\label{alg:greedy}
\begin{algorithmic}[1]
	\State
	\textbf{Input:} $\policy_{\start}, \{\step_{\run}\}_{\run \in \N}, \{\expar_{\run}\}_{\run \in \N}$\vskip 3pt
	\For{$\run = \start,\afterstart,\dots$}\vskip 3pt
	\State
	$\estpolicy_{\run}\leftarrow (1-\expar_\run) \policy_{\run} + \frac{\expar_\run}{\abs{\pures}}$\vskip 3pt
	\State
	Sample $\traj_{\run} \sim \MDP(\estpolicy_{\run}|\rstate_\tstart)$\vskip 3pt
	\State
	$\logvalue_{\run} \leftarrow \reinforce(\estpolicy_{\run},\traj_{\run})$\vskip 3pt
	\State
	$\policy_{\run+1} \leftarrow \proj_{\policies}\parens*{\policy_{\run} + \step_{\run}\logvalue_{\run}}$\vskip 3pt
\EndFor \vskip 5pt
\end{algorithmic}
\end{algorithm}
\end{minipage}
\end{figure*}

%% Reinforce/Greedy algo ends here
%----------------------------------------------------------------------


\begin{model}
[Value-based learning]
\label{mod:value}
The last model we will consider concerns the case where agents only have access to their instantaneous rewards and need to reconstruct their individual gradients based on this information.
A widely used method to achieve this is via the \reinforce subroutine, which we describe in pseudocode form in \cref{alg:reinforce}.
In words, when employing \reinforce, each agent $\play\in\play$ commits to a sampling policy $\hat\policy_{\play}\in\policies_{\play}$ and executes it in an episode of the stochastic game in play.
Then, at the end of the episode, players gather the total reward $\rewards_\play\parens{\traj} \leftarrow \sum_{\time=\tstart}^{\stopTime\parens{\traj}}\reward_{\play,\time}$ associated to the intra-episode trajectory of play $\traj$, and they estimate their policy gradients via the so-called ``log-trick''  \cite{williams1992simple} as
\begin{equation}
\label{eq:logtrick}
\logvalue_\play
	= \rewards_\play\parens{\traj}
		\cdot \sum\nolimits_{\time=\tstart}^{\stopTime\parens{\traj}} \nabla_{\play}
			\parens*{\log\estpolicy_\play\parens{\action_{\play,\time}|\rstate_\time}}.
\end{equation}
\Cref{lem:reinforce} below provides the vital statistics of the \reinforce estimator:

\begin{restatable}{lemma}{ReinforceLemma}
\label{lem:reinforce}
Suppose that each agent $\play\in\players$ follows a stationary policy $\policy_{\play} \in \policies_{\play}$.
Then:
\begin{subequations}
\label{eq:reinforce}
\begin{flalign}
\quad
\text{\itshape a\upshape)}
	\quad
	&\ex_{\traj \sim \MDP}\bracks*{\reinforce(\policy)}
		= \gvalue(\policy)
		&
	\\
\quad
\text{\itshape b\upshape)}
	\quad
	&\ex_{\traj \sim \MDP} \bracks*{\norm{\reinforce_\play(\policy)-\gvalue_\play(\policy)}^2}
		\leq \dfrac{24\nPures_{\play}}{\kappa_\play \stopPr^4}
		&
\end{flalign}
\end{subequations}
where $\kappa_\play = \min_{\rstate\in\states,\pure_{\play}\in\pures_{\play}} \policy_\play(\pure_\play \vert \rstate)$.
\end{restatable}

Therefore, if \reinforce is executed at $\hat\policy \gets \curr$ at each episode $\run=\running$, we will have
\begin{equation}
\exof{\hat\gvalue_{\play,\run}}
%	= \nabla_{\play} \rvalue_{\play,\rstate}(\hat\policy)
	= \gvalue_{\play}(\curr)
	\qquad
	\text{and}
	\qquad
\exof{\dnorm{\noise_{\play,\run}}^{2} \given \curr[\filter]}
	\leq \dfrac{24\nPures_{\play}}{\stopPr^4\min_{\rstate\in\states,\pure_{\play} \in \pures_\play}\policy_{\play,\run}(\pure_\play \vert \rstate)}.
\end{equation}
In particular, this means that we will always have $\curr[\bbound] = 0$ for the bias of the estimator, but its variance could be unbounded if $\curr$ gets close to the boundary of $\policies$.
To avoid this, \reinforce can be paired with an explicit exploration step that modifies the sampling policy of the $\run$-th episode to
\begin{equation}
\estpolicy_{\play,\run}
	= (1-\expar_{\run}) \policy_{\play,\run}
		+ \expar_{\run}\unif_{\pures_{\play}}
		\quad
		\text{\revise{for all $\rstate\in\states$}}
\end{equation}
\ie $\estpolicy_{\play,\run}$ is the mixture between $\policy_{\play,\run}$ and the uniform distribution $\unif_{\pures_{\play}}$ over $\pures_{\play}$.
The resulting algorithm is known as \greedy;
for a pseudocode representation, see \cref{alg:greedy}.

Importantly, by calling \reinforce at $\curr[\estpolicy]$ instead of $\curr[\policy]$, $\curr[\signal]$ becomes biased (because of the difference between $\curr[\estpolicy]$ and $\curr$), but its variance is bounded;
in particular, by invoking \cref{lem:reinforce}, we have
\begin{equation}
\exof{\dnorm{\bias_{\play,\run}} \given \curr[\filter]}
%	\leq \curr[\bbound]
	\leq G\curr[\expar]
	\qquad
	\text{and}
	\qquad
\exof{\dnorm{\noise_{\play,\run}}^{2} \given \curr[\filter]}
%	\leq \curr[\sdev]^{2}
	\leq \frac{24\nPures_{\play}^2}{\expar_{\run}\stopPr^{4}}
\end{equation}
where $G$ is a constant that depends on the smoothness of $\rvalue$ and the cardinalities of $\pures$ and $\states$.%
\footnote{Specifically, from  \cref{lem:smoothness} we know that $\norm{v_i(\hat{\pi}_n) - v_i(\pi_n)} \leq 3\sqrt{\nPures}/\zeta^3 \cdot \sum_{\playalt}\sqrt{\nPures_{\playalt}} \cdot \norm{\hat{\pi}_{j,n}-\pi_{j,n}}$. Moreover,  $|\pi_{i,n}(\pure \mid s)-\hat{\pi}_{i,n}(\pure \mid s)| \leq \varepsilon_n$ for all $\rstate\in\states, \pure \in \mathcal{A}_i$, so $\norm{\pi_{i,n} - \hat{\pi}_{i,n}} \leq \sqrt{\nStates\nPures_{\play}} \varepsilon_n$.
Combining the above, it follows that we can take $G=3\nPlayers\nPures^{3/2}\sqrt{\nStates} \big/ \zeta^3$.}
In this way, \cref{alg:greedy} can be seen as a special case of \eqref{eq:PG} with $\curr[\bbound] = \bigoh(\curr[\expar])$ and $\curr[\sdev]^{2} = \bigoh(1/\curr[\expar])$.
\endenv
\end{model}